\documentclass[10pt,twocolumn]{article}
\usepackage[utf8]{inputenc}
\usepackage[T1]{fontenc}
\usepackage[margin=0.75in,columnsep=0.3in]{geometry}
\usepackage{booktabs}
\usepackage{graphicx}
\usepackage{amsmath}
\usepackage{url}
\usepackage[hidelinks]{hyperref}
\usepackage[numbers,sort&compress]{natbib}

% Revision marks: delete before submission.
\newcommand{\todo}[1]{\textbf{[TODO: #1]}}
% Shows a placeholder when a generated figure has not been built yet.
\newcommand{\genfig}[1]{\IfFileExists{generated/#1.pdf}{\includegraphics[width=\columnwidth]{generated/#1.pdf}}{\fbox{\parbox{0.9\columnwidth}{Run \texttt{python -m nycped paper} to build #1.pdf.}}}}

\title{Exposure-Adjusted Pedestrian Harm Is 16 Percent Higher per 10 Points of Neighborhood Poverty in New York City}
\author{Aditya Bakshi\\\small Independent researcher}
\date{Draft of \today}

\begin{document}
\maketitle

\begin{abstract}
Pedestrian harm is often reported to be higher in poorer neighborhoods, but rates per resident can obscure the contribution of pedestrian exposure. We combine modeled pedestrian volumes for New York City with street centerline data and five years of police-reported crashes (2019 to 2023). We infer crosswalks, group them into 41,937 sites that are mostly intersections, and assign crashes to the nearest site. Among 38,477 sites in tracts with complete data, pedestrian harm crashes rise 16 percent per 10 point increase in the tract poverty rate (rate ratio 1.16, 95 percent interval 1.12 to 1.20) after accounting for crossing volume, street width, travel lanes, posted speed, tract commuting and land use, and borough. Without those adjustments the ratio is 1.30. The estimate holds across estimators, site definitions, time periods, boroughs, alternative standard errors, and a truck route proxy. It falls to 1.10 within ZIP areas and rises to between 1.19 and 1.22 once sampling error in the poverty rate is corrected. Pedestrian deaths show a similar gradient (1.29, 1.17 to 1.42), partly through more crashes and partly through a higher chance that a crash is fatal. A model with residents alone gives ratios of 1.06 to 1.09, so the exposure adjustment, not the choice of geography, drives the result. These are associations from observational data without vehicle traffic volumes. All code is open and rebuilds the numbers from public sources.
\end{abstract}

\section{Introduction}
Between 2019 and 2023, police in New York City recorded 37,815 crashes away from highways in which a pedestrian was hurt or killed, and 469 pedestrians died. These harms fall unevenly across the city, and a common way to show it is to divide crashes by the people who live in each neighborhood.

Residents are a weak denominator. People walk where they work, shop, and change transit, not only where they live. A poorer neighborhood can have a high rate per resident because its streets are more dangerous, because more people cross them, or both. Only the first is a difference in safety. Distinguishing among these explanations requires a measure of how many people walk across each crossing.

Recent modeled pedestrian volumes for New York City \citep{sevtsuk2026nycwalks} make that count possible across the whole city. We use them to measure exposure at 38,477 crossing sites, most of them intersections, and to ask what neighborhood poverty adds once exposure is held fixed.

This paper asks one question. Does neighborhood poverty relate to pedestrian harm after accounting for how many people walk and for street design? We report two disparities. The total gap is the difference that remains with only crossing volume and street design held fixed. The residual gap is the difference that remains after tract commuting and land use are added. The two answer different questions, because transit and walking shares may sit on the causal path from poverty to exposure to harm.

The paper makes three contributions.
\begin{enumerate}
\item It measures exposure at 38,477 crossing sites rather than per resident, and shows that exposure-adjusted pedestrian harm is about 16 percent higher per 10 points of tract poverty (30,866 harm crashes, 411 deaths).
\item It examines the robustness of that estimate to the choice of estimator, the definition of a site, spatial dependence, a truck route proxy for vehicle traffic, sampling error in the poverty rate, error in the modeled volumes, and the choice of geographic scale.
\item It releases a tested package that rebuilds every number from public sources and two small processed inputs.
\end{enumerate}
\section{Related Work}
\paragraph{Exposure in pedestrian safety models.} Pedestrian counts are scarce, so many safety performance functions omit pedestrian exposure or replace it with proxies such as population density or the share of commuters who walk \citep{tao2021exposure}. Tao et al.\ compared Minneapolis models with and without estimated pedestrian and bicycle exposure, entered in logarithmic form alongside vehicle traffic. Including exposure improved fit and changed which locations ranked as highest risk. For pedestrians, only 42 percent of the top 1 percent of intersections were common to both sets of models. For New York City, the NYC Walks authors divided pedestrian injuries at intersections by modeled volume and mapped the result, finding that the intersections with the highest injury risk per pedestrian are often outside Manhattan \citep{sevtsuk2026nycwalks}. That analysis is descriptive and does not relate risk to neighborhood characteristics.

\paragraph{Neighborhood disparities.} The New York City Department of Health reported a pedestrian fatality rate of 9.1 per 100 miles of street in ZIP codes with very high poverty, against 3.2 in low-poverty ZIP codes, for 2012 to 2014, when 497 pedestrians died citywide \citep{nychealth2017}. The denominator is street length, which does not measure how many people walk. The city's current equity reporting under Vision Zero likewise uses rates per resident, comparing killed or severely injured counts per 1,000 population in priority investment areas, which are defined largely by their share of low-income and non-White residents \citep{vzreimagined2026}. In Oregon, rates of fatal and severe pedestrian injury were nearly three times higher in census tracts with the highest concentration of poverty and people of color than in those with the lowest \citep{roll2021oregon}. The Oregon analysis is descriptive. It compares categories and reports higher arterial traffic, more transit stops, and more walking and transit commuting in the high-concentration tracts. In Minneapolis, modeled pedestrian crash risk at intersections was higher in areas of concentrated poverty than elsewhere, with exposure included in the models \citep{lindsey2019equity}. That comparison relies on simple tests of means across areas and on exposure estimated from counts at a few hundred locations.

\paragraph{Contribution.} We combine the two literatures at a larger scale. We estimate the poverty gradient in pedestrian harm at 38,477 crossing sites across a single city, with modeled pedestrian volume at every site, and report how the gradient changes as exposure, street design, and neighborhood context are added. We also test the estimate against sampling error, spatial dependence, a proxy for vehicle traffic, and the choice of geographic scale, which is known to affect area-based results \citep{openshaw1984maup,fotheringham1991mup}. We do not have vehicle volumes, which the Minneapolis models include.

\section{Data}
\subsection{Crashes}
Outcomes come from the NYC Motor Vehicle Collisions Crashes file \citep{nyc_crashes}. We keep crashes from 2019 through 2023 in which at least one pedestrian was injured or killed (pedestrian harm crashes), drop crashes flagged as occurring on highways, and keep those with valid coordinates. This leaves 37,815 crashes, 469 of them fatal to a pedestrian.

\subsection{Pedestrian volumes}
Exposure comes from NYC Walks \citep{sevtsuk2026nycwalks}, which models peak-hour pedestrian volumes on 315,577 sidewalk, crosswalk, and footpath segments. The model is calibrated to New York City Department of Transportation counts from 2018 and 2019. The file gives volumes for six peak hours: weekday and weekend, morning, midday, and evening. We average the six to form one volume index per segment. This is an index of exposure, not a daily count. Because the volumes date from 2018 and 2019 while crashes run to 2023, the index assumes the relative pattern of walking across the city did not change much, which the disruption of 2020 makes questionable. We return to this in the limitations.

\subsection{Streets, tracts, and tract covariates}
Street width, travel lanes, posted speed, road type, and truck route class come from the NYC street centerline file \citep{nyc_centerline}. Tract poverty, population, walking share of commuting, and transit share of commuting come from the 2019 to 2023 American Community Survey five-year estimates \citep{acs2023}. The commercial share of floor area comes from PLUTO \citep{pluto}. Tracts follow the 2020 cartographic boundaries. The tract table has 2,189 tracts, and 2,183 of them contain a site.

\section{Methods}
\subsection{Inferring crosswalks and sites}
The NYC Walks file has no field that identifies crosswalks. We infer them as pedestrian segments that intersect the centerline of an ordinary street, meaning road type 1, active status, and not marked non-vehicular. This flags 110,603 segments, 35 percent of the network. Their median length is 16.6 m, against 64.6 m for all other segments, which is consistent with crosswalks being short crossings and sidewalks and paths being long. The 25 m rule treats crosswalk midpoints within 25 m of each other as one site, using DBSCAN with a minimum of one point. This gives 41,937 sites with a median of three crosswalks per site. Crash coordinates in the source file tend to sit at the intersection center, so a site, not a single crosswalk, is the appropriate unit.

Each site takes the sum of its crosswalk volumes, the widest street it crosses, the highest posted speed, and the most travel lanes. It takes the poverty rate and other covariates of the tract that contains it.

\subsection{Assigning crashes}
A crash belongs to the nearest site within 30 m. About 82 percent of harm crashes meet this rule. Most of the rest are mid-block. The result therefore describes harm near crossings, and we report sensitivity to radii of 20 and 50 m.

\subsection{Models}
For site $i$ in tract $t(i)$ let $Y_i$ be the crash count over the five years. The primary model is Poisson regression with a log link,
\begin{equation}
\log \mathrm{E}[Y_i] = \beta_0 + \beta_P \, P_{t(i)} + \gamma \log V_i + \delta^{\top} S_i + \theta^{\top} Z_{t(i)} + \eta_{b(i)},
\label{eq:model}
\end{equation}
where $P$ is the tract poverty rate in points, $V_i$ is site volume, $S_i$ holds street width in units of 10 ft, travel lanes, and posted speed in units of 5 mph, $Z$ holds crosswalks per site and the standardized walking, transit, and commercial shares, and $\eta_b$ is a borough effect. The quantity of interest is $\exp(10\beta_P)$, the rate ratio per 10 point increase in poverty. Entering log volume as a covariate lets risk rise less or more than proportionally with exposure, which a fixed offset would rule out. The coefficient $\gamma$ is the volume elasticity.

Four specifications show how the ratio changes as controls enter.
\begin{description}
\item[A] Poverty and log volume.
\item[B] Plus street width, lanes, and speed.
\item[C] Plus crosswalks per site and tract controls.
\item[D] Plus borough effects. This is the primary specification.
\end{description}
Poisson estimates stay consistent under overdispersion \citep{santossilva2006log}, so we use them as primary and report standard errors clustered by ZIP area (183 clusters) \citep{cameron2015practitioner}. A negative binomial GLM with a moment-based dispersion parameter is a cross-check. Pedestrian deaths use the same specifications. To split the death gradient, we model harm crashes per site (frequency) and deaths per harm crash among sites with a crash (severity) separately.

\subsection{Robustness and reviewer checks}
We run the following, each with model D unless noted.
\begin{itemize}
\item \emph{Definitions.} Site grouping distance of 15, 25, and 40 m. Crash assignment radius of 20, 30, and 50 m. Dropping the top 1 percent of sites by crash count.
\item \emph{Time and place.} Two periods and each year separately, with the same volumes. Leaving out each borough in turn.
\item \emph{Volume error.} Simulation extrapolation \citep{cook1994simex} at assumed reliabilities of 0.95 to 0.70 for log volume, in the classical-error worst case.
\item \emph{Spatial dependence.} Moran's I \citep{moran1950notes} of residuals using 8 nearest neighbors, standard errors clustered by tract and by 1, 2, and 4 km grid cells, and a block bootstrap over 2 km blocks.
\item \emph{Unmeasured regional factors.} ZIP area fixed effects, which use only variation among tracts in the same ZIP area.
\item \emph{Vehicle traffic.} The data have no vehicle volumes. As a partial proxy we flag sites on a designated truck route and add the flag as a covariate, then split the sample by it.
\item \emph{Sampling error in poverty.} The ACS reports a margin of error for each tract. We treat the margin divided by 1.645 as a known standard error and extrapolate the coefficient back to zero error.
\item \emph{Geographic scale.} The same data aggregated to tracts, ZIP areas, and 100 random contiguous zonings at each of five target sizes \citep{fotheringham1991mup}. Two model forms are compared: residents only (log population plus walking, transit, and commercial shares and borough) and residents plus volume and design.
\end{itemize}

\section{Results}
\subsection{Plausibility of the volumes}
The modeled volumes should track activity, so we compared them with independent tract measures. Median volume per crossing is highest in Manhattan (290) and lower in Brooklyn (162), the Bronx (159), Queens (127), and Staten Island (110). Across tracts the rank correlation of mean volume per crossing is 0.60 with population density and 0.51 with the ACS walking share, and 0.20 with poverty. The spread across boroughs looks narrow relative to density, so volumes may be overstated in low-activity areas. We return to this under measurement error.

\subsection{Pedestrian harm and poverty at crossing sites}
Table~\ref{tab:main} and Figure~\ref{fig:ladder} show the poverty rate ratio as controls enter. Figure~\ref{fig:gradient} shows the same pattern without a poverty term in the model, as observed over expected harm by poverty decile. Sites in the poorest decile (mean tract poverty 38.5 percent) have 3.2 times the harm that crossing volume and street design predict, against sites in the least poor decile (2.7 percent). The ratio is 1.7 once tract context and borough are also held fixed (95 percent intervals 2.65 to 3.89 and 1.44 to 1.97). With only crossing volume held fixed, harm crashes are 30 percent higher per 10 points of poverty (1.30, 1.26 to 1.35). Street design alters the ratio only slightly (1.28). Tract commuting and land use cut it to 1.15, and borough effects leave it at 1.16 (1.12 to 1.20). Together the controls absorb more than 40 percent of the log gap, almost all of it through the tract controls. Street design alone absorbs about 8 percent. The negative binomial cross-check gives the same ratios within 0.06.

\begin{table*}[t]
\centering
\caption{Poverty rate ratio per 10 point increase in tract poverty, with 95 percent intervals. Poisson regression with standard errors clustered by ZIP area. Harm crashes: 38,477 sites, 30,866 crashes. Deaths: 411 deaths.}
\label{tab:main}
\small
\begin{tabular}{lccc}
\toprule
Model & Harm, Poisson & Harm, neg.\ binomial & Deaths, Poisson \\
\midrule
A\quad Volume only & 1.30 (1.26 to 1.35) & 1.35 (1.30 to 1.41) & 1.28 (1.18 to 1.39) \\
B\quad + street design & 1.28 (1.23 to 1.32) & 1.33 (1.27 to 1.39) & 1.26 (1.16 to 1.37) \\
C\quad + tract controls & 1.15 (1.11 to 1.18) & 1.15 (1.11 to 1.19) & 1.23 (1.11 to 1.36) \\
D\quad + borough (primary) & 1.16 (1.12 to 1.20) & 1.17 (1.12 to 1.21) & 1.29 (1.17 to 1.42) \\
\bottomrule
\end{tabular}

\end{table*}

\begin{figure}[t]
\centering
\genfig{ladder}
\caption{Poverty rate ratio per 10 points as controls enter, for harm crashes and deaths. Bars are 95 percent intervals clustered by ZIP area.}
\label{fig:ladder}
\end{figure}

\begin{figure}[t]
\centering
\genfig{gradient}
\caption{Observed over expected pedestrian harm crashes by decile of tract poverty, scaled so the lowest decile is 1 by construction and has no interval. Expected counts come from Poisson models that omit poverty. Bars are 95 percent intervals from resampling ZIP areas.}
\label{fig:gradient}
\end{figure}

In model D, wider streets carry more harm: 26 percent more per additional 10 ft of street width (1.26, 1.21 to 1.31), 9 percent more per travel lane (1.09, 1.04 to 1.15), and 24 percent more per 5 mph of posted speed (1.24, 1.08 to 1.41). Posted speed barely varies in the city, so the speed estimate deserves caution. The volume elasticity is 0.79, meaning harm rises less than proportionally with pedestrian volume. Standardized tract transit share has a large coefficient (1.40 per SD, 1.34 to 1.46), walking share a moderate one (1.11), and commercial share a small negative one (0.93). The poverty gradient does not vary with volume (interaction $p=0.50$) and is slightly smaller on wider streets, by about 3 percent per 10 ft.

\subsection{Robustness}
Figure~\ref{fig:robust} summarizes the checks. The ratio stays between 1.10 and 1.20 in every one that uses the same outcome and covariates. Changing the crash radius to 20 or 50 m gives 1.17 and 1.16. Regrouping crosswalks at 15 or 40 m gives 1.14 and 1.16. Dropping the top 1 percent of sites gives 1.13. The two periods, 2019 to 2021 and 2022 to 2023, both give 1.16, and the yearly ratios run from 1.12 to 1.19. Leaving out each borough in turn gives 1.13 to 1.19, with the lowest value when Manhattan leaves.

\begin{figure}[t]
\centering
\genfig{robustness}
\caption{Poverty rate ratio per 10 points across robustness checks. Bars are 95 percent intervals where the check reports one. Open circles are point estimates. The dotted line marks the primary estimate of 1.16.}
\label{fig:robust}
\end{figure}

\paragraph{Measurement error in volumes.} If the modeled volumes carry classical error with reliability from 0.95 down to 0.70, the extrapolated poverty ratio rises to 1.17 to 1.18 under linear extrapolation and to 1.18 to 1.25 under quadratic extrapolation. The primary estimate is therefore not inflated by noisy volumes. Classical error is a worst case. Model predictions often behave more like Berkson error, which does not bias the coefficient.

\paragraph{Spatial dependence.} Residuals are positively autocorrelated (Moran's I 0.088 at sites, permutation $p=0.005$, and 0.198 when summed by tract), which indicates that some spatially structured factor remains unmodeled. Inference does not depend on it. Clustering by tract, or by 1, 2, or 4 km grid cells, gives intervals within 0.01 of the primary interval, and a block bootstrap over 247 blocks gives 1.16 (1.13 to 1.20).

\paragraph{Differences between ZIP areas.} With ZIP area fixed effects, the ratio is 1.10 (1.07 to 1.13). About a third of the log gap therefore lies between ZIP areas, and a gradient remains among tracts in the same ZIP area. ZIP areas are large, so unmeasured neighborhood factors can still operate inside them.

\paragraph{Vehicle traffic.} Truck route sites carry 1.47 times the harm of other sites after the other controls. Adding the flag leaves the poverty ratio at 1.17 (1.13 to 1.20). The ratio is 1.13 (1.09 to 1.18) off truck routes and 1.20 (1.16 to 1.25) on them. About 17.7 percent of sites qualify. The flag identifies a street class and is not a vehicle count, so this shows only that heavy-vehicle street class does not explain the gradient.

\paragraph{Sampling error in poverty.} The mean standard error of the tract poverty rate is 4.78 points. The reliability of the poverty rate is 0.79, and 0.69 among tracts that are alike on the other controls. Correcting for this error raises the harm ratio from 1.16 to between 1.19 and 1.22 and the death ratio from 1.29 to between 1.34 and 1.41. Sampling error therefore pulls the estimate toward 1 and the primary estimate is conservative. The method assumes independent classical errors. The walking and transit shares are also ACS estimates, so some confounding they would absorb may remain.

\subsection{Pedestrian deaths}
Deaths show a similar and slightly steeper gradient (Table~\ref{tab:main}): 1.29 (1.17 to 1.42) in model D with 411 deaths, 1.28 (1.16 to 1.41) under the negative binomial GLM, and 1.29 (1.16 to 1.42) from a cluster bootstrap over 183 ZIP clusters. Estimates range from 1.23 to 1.35 across site definitions, crash radii, periods, and borough exclusions. The ratio is 1.25 (1.13 to 1.39) for 2019 to 2021 and 1.35 (1.16 to 1.56) for 2022 to 2023, and yearly ratios run from 1.10 to 1.43 with overlapping intervals.

The frequency and severity split gives a harm crash ratio of 1.16 (1.12 to 1.20), a ratio of deaths per harm crash of 1.11 (1.00 to 1.23), and a ratio of deaths per site of 1.29 (1.17 to 1.42). About 60 percent of the log death gradient thus comes from more crashes and about 40 percent from a higher chance that a crash is fatal. The severity interval touches 1, so the second part is suggestive. About 12 percent of deaths (469 against 411) fall outside the 30 m site assignment, so this result describes deaths near crossings.

\subsection{Geographic scale}
\begin{table*}[!t]
\centering
\caption{Poverty rate ratio per 10 points at different zone scales. Poisson regression. ``Residents only'' uses log population, commuting shares, and borough. ``With exposure'' adds log crossing volume and volume-weighted street width, lanes, and speed. Random zones show the median over 100 draws.}
\label{tab:scale}
\small\setlength{\tabcolsep}{4pt}
\begin{tabular}{lcccc}
\toprule
 & \multicolumn{2}{c}{Harm crashes} & \multicolumn{2}{c}{Deaths} \\
\cmidrule(lr){2-3}\cmidrule(lr){4-5}
Scale & Residents only & With exposure & Residents only & With exposure \\
\midrule
Tracts (2,183) & 1.06 (1.02 to 1.10) & 1.13 (1.10 to 1.16) & 1.16 (1.05 to 1.28) & 1.23 (1.12 to 1.36) \\
ZIP areas (183) & 1.09 (1.01 to 1.16) & 1.20 (1.12 to 1.29) & 1.17 (1.00 to 1.36) & 1.25 (1.05 to 1.50) \\
Random, about 2 tracts & 1.07 & 1.15 & 1.15 & 1.24 \\
Random, about 3 tracts & 1.06 & 1.16 & 1.12 & 1.20 \\
Random, about 5 tracts & 1.07 & 1.17 & 1.13 & 1.20 \\
Random, about 9 tracts & 1.07 & 1.18 & 1.11 & 1.19 \\
Random, about 14 tracts & 1.07 & 1.18 & 1.07 & 1.15 \\
\bottomrule
\end{tabular}

\end{table*}

Table~\ref{tab:scale} aggregates the same data to larger zones. With residents only, the poverty ratio for harm crashes is 1.06 to 1.09 at every scale. With volume and design it is 1.13 at tracts, 1.20 at ZIP areas, and 1.15 to 1.18 across random zonings. For a given zone size, 95 percent of random zonings span about plus or minus 0.04. The exposure adjustment, and not the geography, produces the larger estimate. For deaths the exposure-adjusted ratio drifts from 1.29 at sites and 1.23 at tracts to about 1.15 for zones of 14 tracts, and the spread across zonings widens because few deaths fall in each zone.

\section{Discussion}
\paragraph{Principal finding.} After accounting for how many people walk across each crossing and for the design of the street, pedestrian harm crashes are about 16 percent higher per 10 points of tract poverty, and deaths about 29 percent higher. The estimate is stable across the checks we ran and is conservative with respect to sampling error in the poverty rate. Pedestrian exposure and the geographic unit of analysis affect the result in different ways. With residents alone, the ratio is 1.06 to 1.09 at every scale. Adding crossing volume and street design raises it to 1.13 to 1.20. Total crossing volume is negatively correlated with poverty across tracts (Spearman $-0.16$), so poorer tracts carry less pedestrian activity and more harm. A comparison per resident therefore understates the difference in risk at the crossing.

\paragraph{Two disparities.} The total gap (1.28, with volume and street design held fixed) and the residual gap (1.16, after tract context and borough) answer different questions. The total gap describes how much more harm a pedestrian crossing in a poorer tract faces than one with the same volume and street design in a less poor tract. The residual gap asks the same question among tracts that are alike in commuting patterns and land use. Transit share has the largest coefficient in the model (1.40 per standard deviation). It may represent a mediating pathway, if residents of poorer tracts rely more on transit and so walk to and from stops along arterial streets, or it may stand in for density and the character of the street environment. The data cannot separate these readings, so we report both gaps. About a third of the log gap lies between ZIP areas and the rest lies among tracts within them, which suggests that the gradient is not only a difference between broad regions of the city.

\paragraph{Fatal crashes.} The gradient for deaths is similar to that for harm crashes and somewhat steeper. About 60 percent of it reflects more crashes and about 40 percent a higher probability that a crash is fatal, although the second share is imprecise, since its interval includes 1. The data do not allow us to explain a higher fatality rate. Possible explanations include features of the street or vehicle mix that our measures do not capture, and differences in the time taken to receive care. Each would require data we do not have.

\paragraph{A modeling lesson.} A per-resident tract model can miss a gradient that exposure reveals. A tract model of deaths with population as an exact offset and no borough effects gave a poverty ratio of 1.02 (0.92 to 1.12). With population entering freely, borough effects, and the deaths assigned to sites, tract-level ratios are 1.08 to 1.16 without volume and 1.23 with it (Appendix~\ref{app:tract}). The earlier null was fragile to specification and we do not present it as a finding. An offset assumes that risk is proportional to its denominator. The volume elasticity of 0.79 suggests that this assumption does not hold for pedestrian volume. That estimate is itself sensitive to measurement error, however. When we correct for classical error in log volume at assumed reliabilities of 0.95 to 0.70, the elasticity rises from 1.0 to 1.6, so the data do not support a conclusion that risk per pedestrian falls as volume increases. The poverty ratio is much less sensitive to the same correction (1.17 to 1.25).

\paragraph{Limitations.}
\begin{itemize}
\item These are associations. The analysis has no vehicle traffic volumes and uses a truck route flag as a partial proxy. Residual confounding is possible, and residual spatial autocorrelation indicates that a spatial factor remains unmodeled.
\item The pedestrian volumes are modeled, cover six peak hours, and date from 2018 and 2019, while crashes run to 2023. They were calibrated mostly at commercial corridors. Pedestrian activity changed after 2020, although the poverty ratio is stable across periods.
\item About 18 percent of harm crashes and 12 percent of deaths fall outside the 30 m site assignment, mostly mid-block, and are excluded.
\item Tract measures are ecological and say nothing about who is injured.
\item Poverty comes from one ACS period and is measured with error. The transit and walking shares may be mediators, so the residual gap should not be read as the total effect of poverty.
\item Crosswalks are inferred from geometry, not observed.
\end{itemize}

\section{Conclusion}
Pedestrian harm in New York City is higher at crossings in poorer neighborhoods, and the difference persists after accounting for modeled pedestrian volume, street width, travel lanes, posted speed, neighborhood commuting and land use, and borough. At the crossing level the ratio is 1.16 per 10 points of tract poverty for harm crashes and 1.29 for deaths, and it is larger when measurement error in the poverty rate is taken into account. Per-resident comparisons understate it. These findings are associations. They do not identify a cause or the effect of any intervention, and they rest on modeled volumes from 2018 and 2019 and on a proxy for vehicle traffic. Future work should add vehicle volumes and current pedestrian counts, test whether the pattern holds in other cities, and examine which features of the street environment account for the residual gap. The package released with this paper allows each of these steps to be run on the same footing.

\section*{Code and data availability}
All analysis code is in the \texttt{nycped} Python package \citep{nycped_code}, with unit tests, a command line, and a five step reproduction guide. The package rebuilds the site table from three public sources (the NYC Walks supplementary network, the NYC street centerlines, and Census tract boundaries) and two small processed inputs. The processed inputs are a tract table and a table of pedestrian harm crashes, and they are deposited at \todo{Zenodo DOI}. The code is archived at \todo{Zenodo DOI}. The NYC Walks files are not redistributed, since no license is stated. Readers should download them from the source cited in \citet{sevtsuk2026nycwalks}.

\section*{Use of generative AI}
We used Claude, a large language model from Anthropic, to help plan the analysis, write and test code, check results, and draft and edit this text. We reviewed the code and the text and take responsibility for all content, including the results and any errors.

\section*{Acknowledgments}
\todo{Thank the Northeast Big Data Innovation Hub and National Student Data Corps Transportation Data Science Project for the Navigator pathway. Confirm wording with the program. Thank the NYC Walks authors once contacted.}

\bibliographystyle{plainnat}
\bibliography{references}

\appendix
\section{Tract-level models}
\label{app:tract}
Table~\ref{tab:tract} gives the tract-level models, which use the Milestone 5 tract table with crashes counted in tracts. Rate ratios are per 10 points of poverty, from negative binomial regression with standard errors clustered by ZIP area. Models 3 and 4 treat crossing volume as an exposure offset. The data reject proportional risk, so we do not rely on them.

\begin{table*}[t]
\centering
\caption{Tract-level poverty rate ratios per 10 points for pedestrian injuries and deaths.}
\label{tab:tract}
\begin{tabular}{lcc}
\toprule
Model & Injuries & Deaths \\
\midrule
1\quad Residents as exposure, with controls & 1.05 (1.02 to 1.08) & 1.02 (0.92 to 1.12) \\
2\quad Residents as exposure, no controls & 1.12 (1.08 to 1.17) & 1.00 (0.92 to 1.09) \\
3\quad Crossing volume as exposure, with controls & 1.21 (1.16 to 1.25) & 1.17 (1.06 to 1.28) \\
4\quad Crossing volume as exposure, no controls & 1.34 (1.28 to 1.39) & 1.19 (1.10 to 1.30) \\
5\quad Residents and volume as covariates, with controls & 1.11 (1.08 to 1.15) & 1.08 (0.99 to 1.18) \\
6\quad Same as 5, plus borough effects & 1.11 (1.08 to 1.15) & 1.11 (1.01 to 1.22) \\
7\quad Same as 5, Staten Island excluded & 1.12 (1.08 to 1.15) & 1.07 (0.98 to 1.18) \\
\bottomrule
\end{tabular}
\end{table*}
In models 5 and 6 the elasticities for injuries are 0.46 for population and 0.41 for crossing volume.

\end{document}
